\documentclass[10pt,a4paper]{amsart}
\usepackage[margin=19mm]{geometry}
\usepackage{amsmath,amssymb,mathtools}
\usepackage[scr=rsfs,cal=euler]{mathalfa}
\usepackage{xfrac}
\usepackage[unicode,hidelinks,bookmarks=false]{hyperref}
\newcommand{\ie}{i.e.\ }
\newcommand{\cf}{cf.\ }
\newcommand{\resp}{respectively\ }
\newcommand{\Subsection}{\subsection}
\providecommand{\nobreakdash}{\nobreak\hbox{-}}
\setlength{\emergencystretch}{2em}
\multlinegap0pt
\usepackage{bm}
\usepackage{bbm}
\usepackage{dsfont}
\usepackage{xspace}
\usepackage{cancel}
\usepackage{mathtools}
\usepackage{amsthm}
\makeatletter
\@ifundefined{Theorem}{\newtheorem{Theorem}{Theorem}}{}
\@ifundefined{Lemma}{\newtheorem{Lemma}[Theorem]{Lemma}}{}
\makeatother
\title[Tree Buckets and the Reconstruction of Pairs of Phylogenetic]{Tree Buckets and the Reconstruction of Pairs of Phylogenetic Trees}
\author{Sky Basire, Michael Hendriksen}
\date{Source: arXiv:2608.25446v1 (CC BY 4.0)}
\begin{document}
\maketitle

\noindent\emph{Source and licence.} Tree Buckets and the Reconstruction of Pairs of Phylogenetic Trees. Version arXiv:2608.25446v1 (CC BY 4.0), distributed under the Creative Commons Attribution 4.0 International licence, as recorded in the arXiv licence metadata for this exact version and shown on the abstract page https://arxiv.org/abs/2608.25446v1. Full rights notice: licenses/arxiv-2608.25446-CC-BY-4.0.txt.\par\smallskip

\noindent\textit{Editorial scope.} Counted unit: the Theorem whose source label is \texttt{t:pendant.depth.recoverable} in the article (source0.tex lines 445--447, source label \texttt{t:pendant.depth.recoverable}), delivered together with the article's own complete original proof of it (source0.tex lines 449--471, one \texttt{proof} environment of 23 lines). No mathematics is written, repaired or re-derived: statement and proof are the article's own text line for line. Required context retained uncounted: pendant\_depth\_lemma (lines 420--427). Every definition, display and label of those lines is retained unchanged. Omitted from this package and not redistributed: 1--419, 428--444, 448--448, 472--1236. Nothing inside a retained excerpt is omitted or altered: every retained body is delivered exactly as the article's lines are written, so no omission receipt of any kind accompanies this package. Citations: the retained excerpts carry no citation command at all, so no bibliography entry of the article is transcribed and no reference is redistributed. Reference adaptation: none was needed and none was made; every reference inside the delivered unit resolves inside it, so no unresolved reference or citation is produced. Printed slips kept literal: no printed word, symbol, hypothesis or number of the counted statement or of its proof was altered. Layout and class: the article's own class is \texttt{article} with packages including geometry, amsmath, amssymb, amsthm, graphicx, ytableau, booktabs, hyperref, enumerate, enumitem, natbib, tikz, fullpage, tkz-berge; the wrapper is the lane's pinned \texttt{amsart} editorial wrapper at 10pt on A4, which restates the theorem-like environments the retained text needs and adds the article's own macro definitions that the retained text needs (none), with extra wrapper packages (bm, bbm, dsfont, xspace, cancel, mathtools). The delivered numbering is the unit's own. Assets: no retained span contains a figure environment, a graphics-inclusion command, a PDF-page inclusion, a table or a TikZ picture; no image file is copied into this package, the package declares no asset dependency and no image file is redistributed with it. Licence: version arXiv:2608.25446v1 of this article is distributed under the Creative Commons Attribution 4.0 International licence, as recorded in the arXiv licence metadata for this exact version and shown on the abstract page https://arxiv.org/abs/2608.25446v1. A full rights notice is delivered at licenses/arxiv-2608.25446-CC-BY-4.0.txt. Characters: every retained excerpt is pure ASCII, so the delivered engine, XeLaTeX with its default Latin Modern text font, reports no missing character.
\begin{Lemma}[Counting Lemma]\label{pendant_depth_lemma}
    Let $T$ be a rooted binary phylogenetic tree with root balance $1$. Then $T$ has a non-zero pendant depth $k$. Furthermore, either 
    
    \begin{enumerate}
        \item $T$ is a caterpillar tree and the $(n-1)$-multideck consists of $n$ subtrees of pendant depth $n-2$, or
        \item $T$ is not a caterpillar and the $(n-1)$-multideck of $T$ consists of $k$ subtrees of pendant depth $k-1$, and $n-k$ subtrees of pendant depth at least $k$, of which zero, two or four of which are strictly greater than $k$. Furthermore, if there are two or four with pendant depth strictly greater than $k$, they each have the same pendant depth.
    \end{enumerate} 
\end{Lemma}

\begin{Theorem}\label{t:pendant.depth.recoverable}
    Let $T_1$ and $T_2$ be rooted phylogenetic trees, on $n>13$ leaves with pendant depths $k_1$ and $k_2$ respectively. Then the set $\{k_1, k_2\}$ is recoverable from their $(n-1)$-bucket.
\end{Theorem}

\begin{proof}
    We observe that the only trees with a pendant depth of zero that have trees with non-zero pendant depth in their $(n-1)$-multideck are those trees with root balance $2$, as any tree has non-zero pendant depth if and only if it has root balance $1$. In the case of trees with root balance $2$, such a tree will have precisely two trees of non-zero pendant-depth in their $(n-1)$-multideck by Lemma \ref{pendant_depth_lemma}.

    We first consider the case in which the $(n-1)$-bucket contains at least one tree of pendant depth zero. Combining the above observation and Lemma \ref{pendant_depth_lemma}, there are three possibilities for a multideck of a single tree that contains at least one subtree of pendant depth zero - it is a tree with root balance $3$ or greater and all $n$ trees in its $(n-1)$-multideck have pendant depth zero; it is a tree with root balance precisely $2$ and its $(n-1)$-multideck has $n-2$ of pendant depth $0$ and $2$ of pendant depth at least $1$, or it has root balance and pendant depth exactly $1$, so it has one tree of pendant depth $0$ and $(n-1)$ subtrees of pendant depth at least $1$.

    It follows that an $(n-1)$-bucket containing at least one tree with zero pendant depth has one of the following nine options for number of trees of zero pendant depth, denoting root balances by $r_1$ and $r_2$ and supposing $r_1 \le r_2$ without loss of generality: all $2n$ ($r_1,r_2\ge 3$), $2n-2$ ($r_1=2,r_2=3$), $2n-4$ ($r_1=r_2=2$), $n+1$ ($r_1=1,r_2=3$), $n$ ($r_2=3$ and the other tree has pendant depth greater than $1$), $n-1$ ($r_1=1,r_2=2$), $n-2$ ($r_2=2$ and the other tree has pendant depth greater than $1$), $2$ ($r_1=r_2=1$ and both have pendant depth $1$), $1$ ($r_1r_2=1$, one tree has pendant depth $1$, and the other has pendant depth greater than $1$). These are all discrete values when $n \ge 6$, and hence pendant depth is always recoverable for $k_1$ or $k_2 \in \{0,1\}$.

    We can therefore assume that both inputs have pendant depth at least $2$.

    We now consider caterpillar subtrees. Observe that if an input is not a caterpillar, it can have at most four subtrees that are caterpillars by Lemma \ref{pendant_depth_lemma}. Hence, we first check if the highest pendant depth of a subtree is $n-2$. If there are $0$, there are no caterpillars in the input trees and if there are $2n$, both input trees were caterpillars (with pendant depth $(n-1)$). If there are at least $n$, we know one of the input trees are caterpillars (since $n>8$ and two non-caterpillars can contribute at most $8$ caterpillar subtrees to the $(n-1)$-bucket). We can then recover the other pendant depth by considering the smallest pendant depth in the bucket, which must be at least $1$ since we assume input trees have pendant depth at least two, and this smallest pendant depth must be $k-1$ for a pendant depth of $k$ by Lemma \ref{pendant_depth_lemma}. Hence we recover $\{n-1,k\}$.
    
    We can now assume that both input trees have pendant depth at least $2$ and that neither input tree is a caterpillar.

    We now consider the smallest pendant depth of a subtree in the $(n-1)$-bucket. As both trees have pendant depth at least $2$, there is some tree with a minimal pendant depth of $k-1>0$, from which we can immediately identify that some input tree had a pendant depth of $k$. If there are precisely $k$ of them, we know there are two distinct pendant depths for the input trees - if there are more than $k$, we know that the input trees had pendant depths $\{k,k\}$. 

    We then check if there are at least $j$ subtrees with pendant depth $j-1$ for each $j \ge 6$ up to the maximum pendant depth of the subtrees in the $(n-1)$-bucket. If there are, by Lemma \ref{pendant_depth_lemma} $j$ of these can only have come from an input tree with pendant depth of $j$, and so we have recovered $\{j,k\}$ for all cases where at least one is six or greater.

    Hence the only remaining cases are those where the input trees have pendant depth $2 \le j,k \le 5$ and $j \ne k$, and without loss of generality assume $j>k$ (and we always know the value of the smaller pendant depth).

    If $j=k+1$, we simply consider the number of trees with pendant depth $k$. The tree of pendant depth $k$ will contribute $(n-k), (n-k-2)$, or $(n-k-4)$ subtrees of pendant depth $k$, while the tree of pendant depth $k+1$ will contribute $k+1$, for a total of $(n+1), (n-1)$ or $(n-3)$ subtrees of pendant depth $k$. In any case, if $j>k+1$ it would not contribute any subtrees with pendant depth $k$, so this is the only case where the number of subtrees of pendant depth $k$ have a different parity to $n$, and so this case is recoverable.

    Finally, we know that the input tree of pendant depth $j$ must contribute at least $n-j-4>n-9$ (since $j \le 5$) subtrees of pendant depth $j$. Since $n>13$, there will be more than four subtrees, so not all can have been produced by the tree with pendant depth $k<j$ and its none of the pendant depth $k$ subtrees produced by the $k$ input tree since $j>k+1$. Hence we can then immediately recover $j$ too, as it is simply the pendant depth of those (at least) $n-9$ subtrees.
\end{proof}

\end{document}
